\documentclass[11pt,a4paper]{article}

\usepackage{graphicx}
\usepackage{xcolor}
\usepackage[left=2.5cm, right=2.5cm, top=2.5cm, bottom=2.5cm]{geometry}
\usepackage{amsmath,amssymb,amsthm}
\usepackage{fontspec}
\setmainfont{Liberation Serif}
\setsansfont{Liberation Sans}
\setmonofont{Liberation Mono}[Scale=0.86]
\usepackage[russian]{babel}
\usepackage{booktabs}
\usepackage{tabularx}
\newcolumntype{Y}{>{\raggedright\arraybackslash}X}
\usepackage{colortbl}
\usepackage{adjustbox}
\usepackage{enumitem}
\usepackage[htt]{hyphenat}
\usepackage{tcolorbox}
\definecolor{linkblue}{HTML}{1e3a5f}
\definecolor{accgreen}{HTML}{0e7c66}
\definecolor{warnamber}{HTML}{b45309}
\definecolor{rowlight}{HTML}{f0f5fa}
\usepackage[
    colorlinks=true,
    linkcolor=linkblue,
    citecolor=darkgray,
    urlcolor=linkblue,
    bookmarks=true,
    bookmarksnumbered=true,
    unicode=true
]{hyperref}

\widowpenalty=10000
\clubpenalty=10000
\emergencystretch=1.8em
\tolerance=2500
\hbadness=2500
\newcommand{\dd}{\mathrm{d}}

\newtheorem{theorem}{Теорема}[section]
\newtheorem{lemma}[theorem]{Лемма}
\newtheorem{proposition}[theorem]{Предложение}
\newtheorem{corollary}[theorem]{Следствие}
\theoremstyle{definition}
\newtheorem{definition}[theorem]{Определение}
\theoremstyle{remark}
\newtheorem{remark}[theorem]{Замечание}

\newtcolorbox{honestybox}[1][]{colback=warnamber!6, colframe=warnamber!80,
  title={\textbf{Честность:} #1}, fonttitle=\small, boxrule=0.7pt, arc=2pt}
\newtcolorbox{findbox}[1][]{colback=accgreen!7, colframe=accgreen!80,
  title={\textbf{Находка:} #1}, fonttitle=\small, boxrule=0.7pt, arc=2pt}

\hypersetup{
    pdftitle={Задача Чоптюка: постановка, разрешимость, атака на 224 = 4·56 и порог e >= 1},
    pdfauthor={Super Z (Z.ai)},
    pdfsubject={Квартика Клейна, тилинг {7,3}, дефект Кунца, критический коллапс},
    pdfkeywords={Choptuik, Klein quartic, 7-3 tiling, PSL(2,7), Cuntz defect, partial isometry, DSI-4}
}

\begin{document}

\tableofcontents
\newpage

% ==========================================================================
\section{Резюме результатов}\label{sec:summary}
% ==========================================================================

Настоящий отчёт продолжает цикл работ над фреймворком Чоптюка--Исаева
(репозиторий \texttt{choptuik\_\allowbreak ac\_\allowbreak bc}) и отвечает
на три вопроса, поставленных
заказчиком: какова точная постановка задачи Чоптюка и можно ли её решить;
почему наблюдение DSI-4 выделяет знаменатель $224 = 4\cdot 56$ --- константу,
родственную комбинаторике тилинга $\{7,3\}$ квартики Клейна; и можно ли
превратить «лестницу препятствий» дефекта Кунца в теорему, доказав порог
$e \ge 1$ для конструкций, изометричных на носителе. Все вычисления выполнены двумя новыми
детерминированными экспериментами пакета
\texttt{choptuik\_\allowbreak experiments/}, а именно
\texttt{exp6\_\allowbreak census\_\allowbreak 224.py} и
\texttt{exp7\_\allowbreak e1\_\allowbreak threshold.py};
зёрна генераторов зафиксированы, каждый прогон воспроизводится одной командой.

Главные результаты сессии таковы. Во-первых, для класса \emph{носитель-изометрических}
пар --- пар частичных изометрий $X_1, X_2 \in M_n(\mathbb{C})$, в котором лежат
все «структурные» модели фреймворка (перестановочные матрицы, усечённые сдвиги
дерева Кунца) --- доказан порог $e = F(X_1,X_2)/n \ge 1$: значение $5/7$,
минимизирующее дефект по всем матричным парам, в этом классе недостижимо,
поскольку требует сжатия носителя. Вместе с точной формулой покрытия
$F = n + 2\,\mathrm{rank}(I - P - Q)$ и тождеством $F = 2n$ для унитарных пар
это превращает лестницу $5/7 \to 1 \to 2$ в теорему с точной арифметикой
дефекта (раздел~\ref{sec:threshold}). Во-вторых, атака на DSI-4 показала, что
$224 = 4\cdot 56 = 8\cdot 28$ --- \emph{переписная} константа тилинга $\{7,3\}$:
она не менее семи независимыми способами выражается через перепись квартики
Клейна, а дробь $9/224$ --- единственная «переписная» дробь в пределах $0.1\%$
от измеренного $c_{K3}$ (раздел~\ref{sec:attack224}). Главное приобретение ---
родовое семейство $c(g)$: два структурных прочтения $27/672$ совпадают
\emph{только} при $g = 3$, поэтому уже одно измерение амплитуды DSI на любой
старшей поверхности Гурвица различает механизмы. В-третьих, постановка задачи
Чоптюка зафиксирована на двух уровнях (физическом и спектральном), и её
разрешимость оценена на трёх уровнях строгости: численно задача решена,
строгая теория частична, замкнутая форма неизвестна (разделы~\ref{sec:statement}
и~\ref{sec:solvability}).

\begin{table}[htbp]
\centering
\small
\caption{Итоги сессии: результаты и их статус.}
\label{tab:summary}
\begin{tabularx}{\textwidth}{@{}lYY@{}}
\toprule
Результат & Формулировка & Статус \\
\midrule
Теорема~A (\S\ref{sec:threshold}) & $e \ge 1$ для частичных изометрий;
равенство $\Leftrightarrow$ $PQ = 0$, $p + q = n$ & доказана + верифицирована
(0 нарушений на 24\,000 пар) \\
Теорема~B (\S\ref{sec:threshold}) & $F = n + 2\,\mathrm{rank}(I{-}P{-}Q)$
точно; дерево: $F = n + 2$ & доказана (точность $10^{-14}$) \\
Теорема~C (\S\ref{sec:threshold}) & унитарные пары: $F = 2n$ точно &
доказана (точность $10^{-14}$) \\
Лемма~единственности (\S\ref{sec:attack224}) & $9/224$ --- единственная
переписная дробь в $0.1\%$ от $c_{K3}$ & вычислительная, исчерпывающий перебор \\
Теорема о вырождении (\S\ref{sec:attack224}) & $c_I(g) = c_{II}(g)
\Leftrightarrow g = 3$ & доказана (алгебра + сетка $g \le 100$) \\
Дискриминатор (\S\ref{sec:attack224}) & ветви $(g{-}1)^{-1}$ и $(g{-}1)^{-2}$
различимы одним измерением при $g \ge 7$ & фальсифицируемое предсказание \\
\bottomrule
\end{tabularx}
\end{table}

% ==========================================================================
\section{Постановка задачи Чоптюка}\label{sec:statement}
% ==========================================================================

\subsection{Физическая постановка: критический коллапс скалярного поля}

Физическая задача Чоптюка (Чоптюк, 1993) --- задача о гравитационном коллапсе
безмассового скалярного поля в сферической симметрии. Рассматривается система
Эйнштейна---Клейна---Гордона $G_{\mu\nu} = 8\pi\,T_{\mu\nu}[\phi]$ с
$T_{\mu\nu} = \partial_\mu\phi\,\partial_\nu\phi
- \tfrac{1}{2}(\partial\phi)^2 g_{\mu\nu}$ и волновым уравнением
$\Box\phi = 0$; метрика в радиальной калибровке, используемой монографией
(формализм Тьюковского), имеет вид
\begin{equation}\label{eq:metric}
\dd s^2 = -\alpha(r,t)^2\,\dd t^2 + a(r,t)^2\,\dd r^2 + r^2\,\dd\Omega^2,
\end{equation}
а уравнения поля и метрики сводятся к связанной системе для трёх функций
$(\phi, a, \alpha)$: скалярное уравнение
$\tfrac{1}{a}\partial_r(\tfrac{1}{a}\partial_r\phi)
- \tfrac{1}{\alpha}\partial_t(\tfrac{1}{\alpha}\partial_t\phi) = V'(\phi)$
с двумя связями на $\partial_t a$ и $\partial_r\alpha$
(уравнения (9.3)--(9.5) монографии). Семейство гладких начальных данных
параметризуется амплитудой $p$; существует критическое значение $p_*$,
разделяющее дисперсию ($M_{BH} = 0$) и коллапс. Чоптюк обнаружил три
универсальных явления около $p_*$. Первое --- \emph{массовый скейлинг}:
$M_{BH} \propto (p - p_*)^{\gamma}$ с показателем $\gamma \approx 0.374$,
не зависящим от семейства данных. Второе --- \emph{дискретная
самоподобность} (DSS): критическое решение воспроизводит себя по логарифму
масштаба с периодом эха $\Delta \approx 3.44$. Третье --- универсальность
самого критического решения. Механизм объясняется теорией возмущений вокруг
DSS-решения: существует ровно одна неустойчивая мода с показателем роста
$\lambda_0 \approx 2.674$, и $\gamma = 1/\lambda_0$; иллюстрация скейлинга
дана на рис.~\ref{fig:choptuik}.

\begin{figure}[htbp]
\centering
\includegraphics[width=0.98\textwidth]{figs_ru/fig_choptuik.png}
\caption{Слева: универсальный массовый скейлинг $M_{BH} \propto (p-p_*)^\gamma$
для двух различных семейств начальных данных (синтетическая иллюстрация
постановки 1993 года). Справа: измеренный показатель $\gamma \approx 0.374$
против кандидата фреймворка $b_{\mathrm{Ch}} = 1 - \cos(2\pi/7) \approx 0.3765$;
зазор $+0.67\%$.}
\label{fig:choptuik}
\end{figure}

\subsection{«Задача Чоптюка» в монографии: спектральная щель квартики Клейна}

В монографии Исаева «задачей Чоптюка» названа родственная, но иная постановка:
\emph{точное определение спектральной щели} $\Delta = 3.443$ оператора Лапласа
на квартике Клейна
$K\colon x^3y + y^3z + z^3x = 0 \subset \mathbb{CP}^2$ (род $g = 3$,
$\mathrm{Aut}(K) \cong \mathrm{PSL}(2,7)$, $|\mathrm{PSL}(2,7)| = 168$),
снабжённой канонической гиперболической метрикой. Униформизация задаётся
фуксовой группой $\Gamma(2,3,7) = \langle A,B,C \mid A^2 = B^3 = C^7 = ABC = I
\rangle$; спинорные фазы образующих $\delta_A = \pi/2$, $\delta_B = \pi/3$,
$\delta_C = \pi/7$ порождают две поправки к щели, и объединённая формула
монографии
\begin{equation}\label{eq:delta_ch}
\Delta_{\mathrm{Ch}} \;=\; \lambda_1 - \frac{R}{4} + \frac{\delta_C^2}{2}
- \frac{\delta_C^5}{22} \;=\; 3.338 + \frac{\delta_C^2}{2} - \frac{\delta_C^5}{22}
\;\approx\; 3.4379
\end{equation}
воспроизводит наблюдаемое $\Delta = 3.443$ с отклонением $0.15\%$. Здесь же
живёт константа критического показателя фреймворка
\begin{equation}\label{eq:bch}
b_{\mathrm{Ch}}(n) = 1 - \cos\frac{2\pi}{n},
\qquad b_{\mathrm{Ch}}(7) = 0.37651\ldots,
\end{equation}
отождествляемая (с точностью $0.67\%$) с физическим $\gamma \approx 0.374$;
период эха DSS $\Delta \approx 3.44$ дополнительно резонирует со щелью
$3.443$ (расхождение $\sim 0.1\%$). Таким образом, «задача Чоптюка» в широком
смысле --- это программа сведения трёх чисел критического коллапса
($\gamma$, $\Delta$, амплитуды DSI-модуляции $c_{K3} = 0.040177576\ldots$)
к спектру и комбинаторике арифметической кривой рода 3.

\begin{honestybox}[статус связки]
На текущий момент связь квартики Клейна с коллапсом --- аналогия, опирающаяся
на числовые совпадения уровня $0.1\%$--$0.8\%$: щель $\Delta_{\mathrm{Ch}}$
(формула~\eqref{eq:delta_ch}), показатель $b_{\mathrm{Ch}}(7)$ и период эха.
Механизм, выводящий спектр возмущений DSS-метрики из спектра Маасса
арифметической кривой, отсутствует; ниже (раздел~\ref{sec:solvability}) это
зафиксировано как центральный пробел программы.
\end{honestybox}

% ==========================================================================
\section{Разрешима ли задача Чоптюка?}\label{sec:solvability}
% ==========================================================================

Ответ зависит от уровня притязаний; удобно различить три уровня.

\textbf{Уровень I: численное решение.} Здесь задача \emph{решена} и
продолжает уточняться рутинными средствами. Комбинация высокоадаптивных
сеток и теории возмущений вокруг DSS-решения даёт $\gamma = 0.374(1)$,
период эха $\Delta \approx 3.44$ и подтверждает ровно одну неустойчивую
моду; дальнейший прогресс --- вопрос вычислительной инженерии, а не принципа.
Фреймворку этот уровень задаёт жёсткую планку: любое аналитическое
«значение» показателя обязано согласовываться с $\gamma$ до $\sim 10^{-3}$,
иначе оно фальсифицировано. Кандидат $b_{\mathrm{Ch}}(7) = 0.3765$ проходит
эту планку лишь на грани: зазор $+0.67\%$ в пять раз больше современной
неопределённости, и как \emph{точное} утверждение он снят, если только
$\delta$-механика монографии (члены $\delta_C^2/2$, $\delta_C^5/22$ и высшие)
не выведет поправку, приводящую значение к $0.374$. Заметим, что для щели
аналогичная программа выполнена: формула~\eqref{eq:delta_ch} даёт $0.15\%$,
и остаток атрибутирован членам высших порядков голономии.

\textbf{Уровень II: строгая математическая теория.} Здесь картина частичная.
Для скалярного поля в сферической симметрии работы Кристодулу (1991--1999)
устанавливают глобальную картину эволюции: дисперсия, образование чёрных дыр,
инстабильность голых сингулярностей и слабый принцип цензуры. Однако
\emph{дискретная} самоподобность критического решения строгой теории не
поддаётся: существование и единственность DSS-критического решения как
теоремы, доказательство универсальности, строгий счёт неустойчивых мод ---
всё это открытые проблемы математической общей теории относительности.
Техники ограниченной вариации, работающие для CSS-подобных задач, к
осциллирующему DSS-режиму напрямую не переносятся; самоподобная замена
переменных делает коэффициенты уравнений сингулярными по периоду эха. Это
означает, что «решить задачу Чоптюка» в строгом смысле --- программа на
уровне лучших современных задач математической GR, а не одного отчёта.

\textbf{Уровень III: замкнутая форма показателя.} Замкнутая формула для
$\gamma$ неизвестна; более того, нет и содержательного объяснения, почему она
должна существовать: показатель --- это обратный показатель роста
единственного неустойчивого собственного значения линеаризованного оператора
на DSS-фоне, и собственные значения таких операторов, как правило,
трансцендентны. Программа фреймворка утверждает ровно противоположное ---
что $\gamma$ выражается через арифметику рода 3. Чтобы это утверждение стало
теорией, нужны две вещи. Во-первых, \emph{механизм}: явное соответствие между
спектральной задачей DSS-возмущений и объектами квартики (спектром Маасса,
тилингом $\{7,3\}$, спинорными структурами). Во-вторых, \emph{точность}:
объяснение зазора $0.67\%$ между $b_{\mathrm{Ch}}(7)$ и $\gamma$, которое не
сводится к подбору. Наш вклад в эту программу здесь двоякий: теоремы
раздела~\ref{sec:threshold} замыкают целый класс возможных «механизмов»
(изометрические переносы), а родовое семейство
раздела~\ref{sec:attack224} превращает единственное совпадение $27/672$ в
фальсифицируемое однопараметрическое семейство предсказаний --- то есть в
проверяемую физику, а не в нумерологию.

\begin{findbox}[вердикт о разрешимости]
Численно задача Чоптюка решена ($\gamma = 0.374(1)$, $\Delta \approx 3.44$,
одна неустойчивая мода). Строгая теория DSS-критического решения --- крупная
открытая проблема математической GR. Замкнутая форма $\gamma$ через
комбинаторику квартики Клейна --- недоказанная программа фреймворка: её
центральные пробелы --- механизм связи спектров и зазор $0.67\%$; оба пробела
теперь имеют точную формулировку и план атаки.
\end{findbox}

% ==========================================================================
\section{Аналитическая атака на DSI-4: почему $224 = 4\cdot 56$?}\label{sec:attack224}
% ==========================================================================

Наблюдение DSI-4 фиксирует амплитуду лог-периодической модуляции $c_{K3}$,
расширенная запись которой равна $0.04017757639214903$, а её целочисленное
прочтение даёт $c_{K3} \approx 27/672 = 9/224$ (отклонение $+0.0025\%$).
Вопрос настоящего раздела: выводит ли комбинаторика тилинга $\{7,3\}$
знаменатель $224 = 4\cdot 56$ естественно и единственно. Ответ: да ---
константа \emph{переписная}, и её переписной статус теперь доказан
конструктивно.

\subsection{Перепись тилинга $\{7,3\}$ из группы $G = \mathrm{GL}(3,2)$}

Мы построили комбинаторику тилинга явно, из группы
$G = \mathrm{GL}(3,2) \cong \mathrm{PSL}(2,7)$ из 168 матриц $3\times 3$ над
$\mathbb{F}_2$. Перепись элементов по порядкам даёт $1, 21, 56, 42, 48$
элементов порядков $1, 2, 3, 4, 7$ соответственно; силовские числа равны
$n_2 = 21$, $n_3 = 28$, $n_7 = 8$. Ячейки тилинга --- классы смежности по
стабилизаторам: вершины $G/H_3$ ($H_3 \cong C_3$, три семиугольника в каждой
вершине), рёбра $G/H_2$, грани $G/H_7$; это даёт $V = 56$, $E = 84$,
$F = 24$ и эйлерову характеристику $V - E + F = -4 = 2 - 2g$. Целостность
построения проверена двумя нетривиальными графами: 1-скелет тилинга
(граф Клейна) --- связный кубический граф на 56 вершинах с 84 рёбрами и
обхватом 7 (единственность такого графа известна), а граф граневой смежности
--- 7-валентный граф на 24 вершинах с законными треугольниками (три
семиугольника в каждой вершине попарно смежны). Три потока валентностей
совпадают: $7F = 3V = 2E = |G| = 168$ --- тождество регулярного тилинга
$\{7,3\}$ рода 3.

\begin{table}[htbp]
\centering
\small
\caption{Перепись квартики Клейна и тилинга $\{7,3\}$ (все значения получены
явным построением из $\mathrm{GL}(3,2)$ и проверены программно).}
\label{tab:census}
\begin{tabularx}{\textwidth}{@{}lYY@{}}
\toprule
Объект & Формула / происхождение & Значение \\
\midrule
Порядок группы & $|\mathrm{PSL}(2,7)| = 84(g-1)$ & 168 \\
Элементы порядков 2 / 3 / 4 / 7 & перепись $\mathrm{GL}(3,2)$ & 21 / 56 / 42 / 48 \\
Силовские числа $n_2, n_3, n_7$ & подгруппы $C_2, C_3, C_7$ & 21 / 28 / 8 \\
Вершины / рёбра / грани & $|G|/3,\ |G|/2,\ |G|/7$ & 56 / 84 / 24 \\
Флаги (ориентированные / все) & $|G|,\ 2|G|$ & 168 / 336 \\
1-скелет $\{7,3\}$ (граф Клейна) & двойной косет $H_3\,t\,H_3$, $t^2 = e$ &
56 вершин, ст.~3, обхват 7 \\
Граневая смежность $\{7,3\}$ & двойной косет $H_7\,t\,H_7$ & 24 вершины,
ст.~7, 84 ребра \\
Эйлер & $V - E + F = 2 - 2g$ & $-4$ \\
Тождество валентностей & $7F = 3V = 2E = |G|$ & 168 \\
\bottomrule
\end{tabularx}
\end{table}

\subsection{Провенанс констант 224 и 672}

Оба знаменателя наблюдения DSI-4 переписываются через перепись
табл.~\ref{tab:census} многократно и независимо; перечисление точных тождеств
даёт критерий «переписности»: константа не случайна, если она лежит в
пересечении многих независимых счётных семейств. Для $224$ таких появлений
семь, для $672$ --- восемь; особо отметим два межбашенных: $224 =
14(g{-}1)|_{g=17}$ --- число $C_3$-подгрупп поверхности триплета Гурвица
рода 17, а $672 = 4V|_{g=7}$ --- учётверённое число вершин поверхности
Фрикке---Макбита рода 7.

\begin{table}[htbp]
\centering
\small
\caption{Провенанс: все точные появления $224$ и $672$ в переписи
$\{7,3\}$ и башни Гурвица.}
\label{tab:prov}
\resizebox{\textwidth}{!}{%
\begin{tabular}{ll@{\hspace{2.2em}}ll}
\toprule
\multicolumn{2}{c}{$224$} & \multicolumn{2}{c}{$672$} \\
\cmidrule(r){1-2}\cmidrule(l){3-4}
$4\cdot 56$ & вершины $\{7,3\}$: $4V$ & $4\cdot 168$ & $4|G|$ (4 спинорных компоненты) \\
$8\cdot 28$ & $8\,n_3$ ($C_3$-подгруппы) & $2\cdot 336$ & $2|\mathrm{PGL}(2,7)|$ \\
$8\cdot 28$ & 28 битангентов $= 8N$ & $24\cdot 28$ & $F\cdot n_3$ (грани $\times$ $C_3$) \\
$32\cdot 7$ & $2^5\cdot 7$ (7 = порядок $C$) & $12\cdot 56$ & $12V$ \\
$168+56$ & $|G| + V$ & $8\cdot 84$ & $8E$ \\
$2\cdot 84+56$ & $2E + V$ & $3\cdot 224$ & 3 = валентность вершины \\
$14\cdot 16$ & $N(g{=}17)$: триплет Гурвица & $4\cdot 168$ & $4V$ рода 7 (Фрикке--Макбит) \\
\bottomrule
\end{tabular}%
}
\end{table}

\subsection{Лемма единственности: $9/224$ среди переписных дробей}

Переписной полугруппой назовём числа $q = 2^a 3^b 7^c$ с бюджетом показателей
$a \le 6$, $b \le 1$, $c \le 1$ --- ровно таковы все знаменатели переписи
(64 спиновых структур, 28 битангентов, 24 грани, 56 вершин, 84 ребра, 168,
336, 672); знаменатели вроде $25 = 5^2$, $49 = 7^2$, $81 = 3^4$, $11^2$
в переписи не встречаются \emph{по построению}. Исчерпывающий перебор дробей
$p/q$ с $q \le 5000$ из полугруппы и $p \le 32$ даёт: в пределах $0.1\%$ от
измеренного $c_{K3}$ лежит \emph{ровно одна} рациональная дробь ---
$9/224$ (её несокращённые записи $18/448$, $27/672$ исчерпывают все попадания;
в пределах $0.01\%$ --- та же единственная дробь). Цепная дробь измеренного
значения $[0; 24, 1, 8, 19, \ldots]$ имеет ровно две переписные подходящие
дроби: $1/24$ (отклонение $+3.68\%$) и $9/224$ ($+0.0025\%$); промежуточная
$1/25$ отбрасывается фильтром как непереписная. Тем самым исключительность
$9/224$ --- не артефакт подгонки: любая альтернатива с сопоставимой точностью
обязана иметь знаменатель вне переписи $\{7,3\}$ (рис.~\ref{fig:scan}).

\begin{figure}[htbp]
\centering
\includegraphics[width=0.94\textwidth]{figs_ru/fig224_scan.png}
\caption{Переписной фильтр: ошибки $|p/q - c_{K3}|$ для всех переписных
знаменателей (синие) и случайной выборки прочих дробей (серые). В полосе
$0.1\%$ --- только $9/224 = 27/672$.}
\label{fig:scan}
\end{figure}

\subsection{Родовое семейство $c(g)$ и теорема о вырождении}

Главный результат атаки: наблюдение $27/672$ допускает \emph{два} структурных
прочтения, которые на квартике Клейна неразличимы, но расходятся на всей
остальной башне Гурвица. Пусть поверхность башни имеет род $g$, порядок
$|G| = 84(g-1)$ и канонический тилинг $\{7,3\}$: тогда $V = 28(g-1)$,
$F = 12(g-1)$ и $N := n_3 = 14(g-1)$ --- число $C_3$-подгрупп (вершинных
стабилизаторов). Ветвь~I читает $27/672$ как $3^3/(4|G|)$ --- «резонансы на
спинорные слоты по группе»; ветвь~II --- как $(N-1)/(FN)$ --- «все, кроме
одной, $C_3$-подгруппы на грань$\times$подгруппу слотов» (прочтение
$(N{-}1)/(24N)$ из эксперимента 2). Явно:
\begin{equation}\label{eq:branches}
c_I(g) = \frac{27}{4\cdot 84(g-1)} = \frac{9}{112(g-1)},
\qquad
c_{II}(g) = \frac{14(g-1) - 1}{12(g-1)\cdot 14(g-1)}.
\end{equation}

\begin{theorem}[о вырождении]\label{thm:degen}
Ветви $c_I(g)$ и $c_{II}(g)$ совпадают тогда и только тогда, когда $g = 3$.
\end{theorem}
\begin{proof}
Равенство $27\cdot 168(g-1)^2 = 336(g-1)\bigl(14(g-1)-1\bigr)$ после
сокращения на $168(g-1)$ даёт $27(g-1) = 28(g-1) - 2$, то есть $g - 1 = 2$.
\end{proof}

Численная проверка на сетке $g = 3..100$ подтверждает: вырожденная точка
единственна. Ветвь~III --- гипотеза фреймворка $c_{III}(g) = b_{\mathrm{Ch}}(b_2(g))$
с гипотетическими $b_2 = 22, 46, 94, 190$ --- отделяется от обеих уже при
$g = 7$ и ведёт себя качественно иначе: $b_{\mathrm{Ch}}(\lambda)
= 2\sin^2(\pi/\lambda) \sim 2\pi^2/\lambda^2$, то есть $(g-1)^{-2}$ против
$(g-1)^{-1}$. Итоговая таблица предсказаний:

\begin{table}[htbp]
\centering
\small
\caption{Родовое семейство амплитуды DSI $c(g)$: три ветви на башне Гурвица.
При $g = 3$ все ветви сливаются (вырождение квартики); при $g \ge 7$ ветви
различимы.}
\label{tab:genus}
\begin{tabular}{rrrrrrr}
\toprule
$g$ & $N = n_3$ & $c_I$ & $c_{II}$ & $c_{II}/c_I$ & $c_{III} = b_{\mathrm{Ch}}(b_2)$ &
$c_{III}/c_I$ \\
\midrule
3 & 28 & $9/224 = 0.040179$ & $9/224 = 0.040179$ & 1.0000 & 0.040509 & 1.0082 \\
7 & 84 & $3/224 = 0.013393$ & $83/6048 = 0.013724$ & 1.0247 & 0.009314 & 0.6954 \\
14 & 182 & $9/1456 = 0.006181$ & $181/28392 = 0.006375$ & 1.0313 & 0.002233 & 0.3613 \\
17 & 224 & $9/1792 = 0.005022$ & $223/43008 = 0.005185$ & 1.0324 & 0.000547 & 0.1089 \\
\bottomrule
\end{tabular}
\end{table}

\begin{findbox}[дискриминатор]
Ветвь~I убывает как $(g-1)^{-1}$, ветвь~III --- как $(g-1)^{-2}$: при
$g = 17$ они различаются в девять раз. \emph{Одного} измерения амплитуды
DSI-модуляции на любой старшей поверхности Гурвица (ближайшие кандидаты:
Фрикке--Макбит $g = 7$ с $c_I = 3/224$, или триплет $g = 17$) достаточно,
чтобы выбрать механизм. Тем самым DSI-4 повышает статус: от единственного
совпадения --- к однопараметрическому фальсифицируемому семейству
(рис.~\ref{fig:genus}).
\end{findbox}

\begin{figure}[htbp]
\centering
\includegraphics[width=0.94\textwidth]{figs_ru/fig224_genus.png}
\caption{Родовое семейство $c(g)$: ветвь~I $9/(4V)$, ветвь~II $(N-1)/(FN)$
и ветвь~III $b_{\mathrm{Ch}}(b_2)$. Точки вырождения $c_I = c_{II}$ ---
ровно $g = 3$ (теорема~\ref{thm:degen}).}
\label{fig:genus}
\end{figure}

\subsection{Статус DSI-4 после атаки}

Итог раздела таков. Знаменатель $224 = 4\cdot 56$ --- переписная константа
тилинга $\{7,3\}$: он выражается через перепись квартики семью способами и
обладает межбашенными двойниками ($N$ триплета рода 17, $4V$ рода 7), что
делает случайным его появление крайне маловероятным. Дробь $9/224$ ---
единственная переписная дробь в окрестности измеренного значения (лемма
единственности, исчерпывающий перебор). Числитель $9 = 3^2$ при этом получает
две интерпретации --- «$3^3$ резонансов на $4|G|$ спинорных слотов» (ветвь~I)
и «$N - 1$ из $FN$ слотов» (ветвь~II) --- и различить их способна только
экспериментальная башня: на роде 3 механизмы вырождены теоремой~\ref{thm:degen}.
Что осталось открытым: динамический вывод числителя (почему именно «все,
кроме одного» или именно $3^3$) --- это та же задача о механизме, что и для
связки с $\gamma$; её постановка теперь точна.

% ==========================================================================
\section{Порог $e \ge 1$: лестница становится теоремой}\label{sec:threshold}
% ==========================================================================

\subsection{Класс носитель-изометрических пар}

Функционал дефекта Кунца для пары $X_1, X_2 \in M_n(\mathbb{C})$:
\begin{equation}\label{eq:F}
F(X_1, X_2) = \|X_1^*X_1 - I\|^2 + \|X_2^*X_2 - I\|^2
+ \|X_1^*X_2\|^2 + \|X_1X_1^* + X_2X_2^* - I\|^2,
\qquad e = F/n,
\end{equation}
где нормы фробениусовы; точные соотношения Кунца $X_i^*X_j = \delta_{ij}I$,
$X_1X_1^* + X_2X_2^* = I$ не реализуются ни в одной $M_n$. Острая теорема
следа аудита даёт глобальный пол $F \ge 5n/7$ с равенством на паре
$\sqrt{4/7}\,I$ --- конструкции, \emph{сжимающей} всё пространство. Однако
естественные «переносы» фреймворка не сжимают носитель. Формализуем это.

\begin{definition}[носитель-изометрическая пара]\label{def:carrier}
$X \in M_n(\mathbb{C})$ --- \emph{частичная изометрия}, если все её ненулевые
сингулярные числа равны 1 (эквивалентно: $XX^*X = X$; носитель --- начальное
подпространство $(X^*X)\mathbb{C}^n$, и на нём $X$ изометрична). Пара
$(X_1, X_2)$ --- \emph{носитель-изометрическая}, если оба оператора ---
частичные изометрии.
\end{definition}

В классе определения~\ref{def:carrier} лежат все структурные модели фреймворка:
перестановочные матрицы (унитарны), усечённые сдвиги дерева Кунца
(частичные изометрии с проекционными $P_i = S_iS_i^*$), комплементарные
проекции. Вне класса --- $\sqrt{4/7}\,I$ (сингулярные числа $\sqrt{4/7}
\approx 0.756$) и взвешенные деревья $c_iS_i$ с $c_i \ne 1$.

\subsection{Теорема A: порог}

\begin{theorem}[порог $e \ge 1$]\label{thm:A}
Для любой носитель-изометрической пары $(X_1, X_2) \in M_n(\mathbb{C})^2$:
\begin{equation}\label{eq:thmA}
F(X_1, X_2) \ge n, \qquad \text{т.е. } e \ge 1.
\end{equation}
Равенство достигается тогда и только тогда, когда $P = X_1X_1^*$ и
$Q = X_2X_2^*$ ортогональны ($PQ = 0$) и $p + q = n$, где $p = \mathrm{rank}\,P$,
$q = \mathrm{rank}\,Q$.
\end{theorem}

\begin{proof}
Если $X_i$ --- частичные изометрии, то $P = X_1X_1^*$, $Q = X_2X_2^*$ ---
ортогональные проекции. Тождество аудита $F = G(P,Q)$ с
$G(P,Q) = 2\,\mathrm{tr}P^2 + 2\,\mathrm{tr}Q^2 + 3\,\mathrm{tr}(PQ)
- 4\,\mathrm{tr}P - 4\,\mathrm{tr}Q + 3n$ вместе с $\mathrm{tr}(P^2) =
\mathrm{tr}(P)$ для проекций даёт
\begin{equation}\label{eq:Fproj}
F = 3n - 2p - 2q + 3r, \qquad r = \mathrm{tr}(PQ).
\end{equation}
Для проекций выполнены оценки
$\max(0,\, p + q - n) \le r \le \min(p, q)$:
нижняя следует из $\mathrm{tr}(PQ) \ge \mathrm{tr}(P \wedge Q) \ge p + q - n$
(размерность пересечения образов), верхняя --- из $\mathrm{tr}(PQ) =
\mathrm{tr}(QPQ) \le \|P\|\,\mathrm{tr}(Q) = q$. Если $p + q \ge n$, то
$F - n = 3r - 2(p+q-n) \ge 3(p+q-n) - 2(p+q-n) = p+q-n \ge 0$; если
$p + q \le n$, то $F - n = 2(n - p - q) + 3r \ge 0$. В обоих случаях
равенство требует $p + q = n$ и $r = 0$, то есть $PQ = 0$.
\end{proof}

\begin{remark}[геометрия равенства]\label{rem:equality}
Случай равенства --- пара комплементарных ортогональных проекций
(например, $X_1 = \mathrm{diag}(I_p, 0)$, $X_2 = \mathrm{diag}(0, I_q)$,
$p + q = n$): дефект $F = n$ целиком сидит в членах «полноты»
$\|X_i^*X_i - I\|^2 = q$ и $p$, а члены перекрытия и покрытия нулевые.
Такая пара вырожденна как модель Кунца: операторы не изометричны на всём
пространстве. Групп калибровки $F$ велик --- $(X_1, X_2) \mapsto
(UX_1W,\, UX_2V)$ с произвольными унитарными $U, V, W$, --- и все случаи
равенства эквивалентны указанному блочному виду.
\end{remark}

\subsection{Теорема B: точная арифметика покрытия}

\begin{theorem}[формула покрытия]\label{thm:B}
Пусть $X_1, X_2$ --- частичные изометрии и $P + Q \le I$ (комбинированные
образы не накрывают пространство), $R := I - P - Q \ge 0$. Тогда
\begin{equation}\label{eq:thmB}
F = n + 2\,\mathrm{rank}(R) \quad \text{точно.}
\end{equation}
\end{theorem}
\begin{proof}
Из $P + Q \le I$ для проекций следует $PQ = 0$. Прямой подсчёт четырёх
слагаемых~\eqref{eq:F} даёт: $\|X_1^*X_1 - I\|^2 = \|P' - I\|^2$, где
$P' = X_1^*X_1$ --- начальная проекция ранга $p$; так как $I - P = Q + R$,
то $\|P' - I\|^2 = \mathrm{tr}(I - P') = (n - p) = q + \mathrm{rank}(R)$;
симметрично $\|X_2^*X_2 - I\|^2 = p + \mathrm{rank}(R)$; далее
$\|X_1^*X_2\|^2 = \mathrm{tr}(PQ) = 0$ и $\|P + Q - I\|^2 = \|R\|^2 =
\mathrm{rank}(R)$. Суммируя и используя $p + q + \mathrm{rank}(R) = n$,
получаем $F = n + 2\,\mathrm{rank}(R)$.
\end{proof}

Для усечённого дерева Кунца глубины $k$ непокрытым остаётся только корень:
$\mathrm{rank}(R) = 1$, и формула~\eqref{eq:thmB} мгновенно воспроизводит
эмпирическую находку эксперимента 1: $F = n + 2$, $e = 1 + 2/n \to 1$
(рис.~\ref{fig:tree}; проверено до $k = 10$, $n = 2047$, точность $10^{-13}$).
Обрезка дерева --- \emph{почти-экстремаль} класса: платит ровно $2$ за
единственную непокрытую размерность.

\begin{figure}[htbp]
\centering
\includegraphics[width=0.88\textwidth]{figs_ru/fig7_tree.png}
\caption{Арифметика покрытия: обрезка дерева Кунца точно следует
$F = n + 2$ (теорема~\ref{thm:B}, $\mathrm{rank}(R) = 1$); пол класса
$F = n$ недостижим для дерева, унитарные пары платят $2n$.}
\label{fig:tree}
\end{figure}

\subsection{Теорема C и трёхчленная арифметика дефекта}

\begin{theorem}[унитарные пары]\label{thm:C}
Если $X_1, X_2$ унитарны, то $F = 2n$ точно, при любой паре.
\end{theorem}
\begin{proof}
$P = Q = I$, и первые два слагаемых~\eqref{eq:F} нулевые,
$\|X_1^*X_2\|^2 = \mathrm{tr}(I) = n$ (произведение унитарно), последнее
слагаемое $\|I\|^2 = n$.
\end{proof}

Перестановочные модели фреймворка --- частный случай: любая унитарная пара
платит ровно $2n$, и дефект целиком сидит в \emph{перекрытии} образов
($r = n$). Формулы~\eqref{eq:Fproj} и~\eqref{eq:thmB} дают точную
трёхчленную арифметику: дефект носитель-изометрической пары складывается из
(i) \emph{перекрытия} образов с весом $3r$, (ii) \emph{недопокрытия} с весом
$2\,\mathrm{rank}(R)$ и (iii) базовой единицы $n$; лестница фреймворка ---
это ровно значения этой арифметики на трёх семействах (рис.~\ref{fig:ladder}):
свободные скаляры $5/7$ (вне класса), дерево $1 + 2/n$, унитарные пары $2$.

\begin{figure}[htbp]
\centering
\includegraphics[width=0.94\textwidth]{figs_ru/fig7_ladder.png}
\caption{Лестница как теорема: на классе носитель-изометрических пар пол
$e = 1$ (теорема~\ref{thm:A}); значение $5/7$ достижимо только сжатием
носителя (вне класса), унитарные пары платят ровно $e = 2$
(теорема~\ref{thm:C}).}
\label{fig:ladder}
\end{figure}

\subsection{Следствие для цепочки Ш1--Ш4 и моста Ш.3(iii)}

\begin{corollary}[замыкание класса]\label{cor:D}
На классе носитель-изометрических пар $\inf e = 1$; значение $5/7 < 1$
недостижимо. Следовательно, никакая последовательность носитель-изометрических
моделей не даёт $e \to 0$: в цепочке Ш1--Ш4 шаг Ш2 усиливается до
$\eta \equiv 1$, и мост Ш.3(iii) (софичность $\Rightarrow$ матричные модели
с $e \to 0$) невозможен для изометрических на носителе переносов --- какие бы
перестановочные данные ни брать, дефект на душу не опускается ниже 1.
\end{corollary}

\begin{remark}[граница класса оптимальна]
Взвешенное дерево $c_iS_i$ покидает класс с квадратичным штрафом
\begin{equation}\label{eq:weightedtree}
F = (n-1)[(c_1^2{-}1)^2 + (c_2^2{-}1)^2] + n + 2;
\end{equation}
при $c^2 = 4/7$
получается $e \approx 1.37$--$1.48$ (для $k = 3$--$7$) --- выше пола класса.
Чтобы добраться до $5/7$, нужно сжать \emph{весь} носитель
($\sqrt{4/7}\,I$): глобальный минимум живёт строго вне класса, и потому
порог теоремы~\ref{thm:A} неустраним внутри класса.
\end{remark}

\subsection{Численная верификация}

Все утверждения проверены экспериментом \texttt{exp7\_e1\_threshold.py}
(детерминированный, зафиксированные зерна); итог --- в
табл.~\ref{tab:verify}. Отдельно проверены опорные оценки
$\max(0, p+q-n) \le \mathrm{tr}(PQ) \le \min(p,q)$ на $10^4$ случайных пар
проекций (0 нарушений) и тождество~\eqref{eq:Fproj} с точностью
$5.9\cdot 10^{-16}$.

\begin{table}[htbp]
\centering
\small
\caption{Численная верификация теорем A/B/C (эксперимент 7).}
\label{tab:verify}
\begin{tabularx}{\textwidth}{@{}lYY@{}}
\toprule
Проверка & Объём & Макс. отклонение \\
\midrule
Тождество $F = 3n - 2p - 2q + 3r$ & 24\,000 случайных частичных изометрий,
$n \le 12$ & $5.9\cdot 10^{-16}$ \\
Порог $e \ge 1$ (теорема~A) & те же 24\,000 пар & 0 нарушений \\
Случай равенства $F = n$ & 5 конфигураций $(n, p, q)$ & $8.9\cdot 10^{-16}$ \\
Формула покрытия (теорема~B) & 500 случайных $(P, Q)$, $P + Q \le I$ &
$7.1\cdot 10^{-15}$ \\
Дерево: $F = n + 2$, $\mathrm{rank}(R) = 1$ & $k = 2..10$ ($n$ до 2047) &
$2.3\cdot 10^{-13}$ \\
Унитарные пары $F = 2n$ (теорема~C) & 200 пар, $n \le 32$ & $2.8\cdot 10^{-14}$ \\
Перестановочные $F = 2n$ & 200 пар & $1.4\cdot 10^{-14}$ \\
Оценки $\mathrm{tr}(PQ)$ & 10\,000 пар проекций & 0 нарушений \\
\bottomrule
\end{tabularx}
\end{table}

% ==========================================================================
\section{Выводы и следующие шаги}\label{sec:conclusions}
% ==========================================================================

Сведение результатов этой сессии в одну картину даёт три сдвига в программе.
Первый: «лестница препятствий» дефекта Кунца перестала быть эмпирической
последовательностью значений и стала теоремой с точной арифметикой --- порог
$e \ge 1$ на классе носитель-изометрических конструкций доказан
(теорема~\ref{thm:A}), формула покрытия $F = n + 2\,\mathrm{rank}(R)$ точна
(теорема~\ref{thm:B}), унитарные пары платят ровно $2n$
(теорема~\ref{thm:C}), и следствие~\ref{cor:D} закрывает мост Ш.3(iii) для
всего класса: софическое усиление не может приходить от изометрических
переносов. Второй: наблюдение DSI-4 получило аналитическое окружение ---
константа $224 = 4\cdot 56$ переписна (семь независимых тождеств, межбашенные
двойники), дробь $9/224$ единственна среди переписных (лемма
единственности), а родовое семейство $c(g)$ с единственным вырождением при
$g = 3$ (теорема~\ref{thm:degen}) делает программу фальсифицируемой: ветви
$(g-1)^{-1}$ и $(g-1)^{-2}$ различаются уже при $g = 7$ в полтора раза и в
девять раз при $g = 17$. Третий: постановка задачи Чоптюка зафиксирована на
обоих уровнях (физическом и спектральном), и её разрешимость оценена честно:
численно --- решено; строго --- крупные открытые проблемы (существование,
единственность, устойчивость DSS); замкнутая форма --- программа фреймворка
с двумя точными пробелами (механизм и зазор $0.67\%$).

Естественные следующие шаги таковы. (а)~Внести теоремы~\ref{thm:A}--\ref{thm:C}
и следствие~\ref{cor:D} в репозиторий как версию 3 \texttt{stability\_lemma}
с тестами на частичных изометриях. (б)~Поставить решёточный/спектральный
эксперимент на поверхности Фрикке--Макбита ($g = 7$): измерение амплитуды
DSI-модуляции различило бы ветви~I и~III уже одним числом ($3/224$ против
$b_{\mathrm{Ch}}(46)$ --- расхождение в $1.44$ раза). (в)~Высокоточный расчёт
$\lambda_1$ квартики Клейна (алгоритм Хейхала для маассовых форм) и показателя
$\gamma$ методом возмущений Гундлаха --- чтобы проверить, действительно ли
$\delta$-разложение монографии способно закрыть зазор $0.67\%$ между
$b_{\mathrm{Ch}}(7)$ и $\gamma$. (г)~Аналитика числителя $9 = 3^2$: попытка
вывода «все, кроме одной» ($N - 1$) из когомологий спинорных пучков на
$\{7,3\}$ либо доказательство отсутствия такого вывода --- что само по себе
усилило бы честную атрибуцию DSI-4 как наблюдения.

\end{document}
